\section{Extended Design and Blueprint Notes}

This section preserves the design-space, anti-pattern, blueprint, and case-study detail that the main paper compresses to stay within the CSUR page budget. It is organized to mirror the main manuscript's argumentative flow: recurring principles, architecture axes, contract repairs, and concrete boundary examples.

\subsection{Design Guide Summary}

The summary below is intentionally short: it turns the comparative synthesis into a compact design checklist. The table names recurring design questions; the following subsections unpack them into principles, axes, and repair patterns.

\begin{table}[t]
\small
\centering
\caption{Extended reader guide for translating synthesis into runtime design choices.}
\label{tab:supp-design-guide}
\begin{tabularx}{\linewidth}{>{\raggedright\arraybackslash}p{0.25\textwidth}Y>{\raggedright\arraybackslash}p{0.22\textwidth}}
\toprule
Recurring synthesis result & Design question & Concrete artifact to look for \\
\midrule
Observability gaps dominate failure analysis & Which signal names the state object, lifecycle stage, delay semantics, and legal action? & typed telemetry and disturbance indicators \\
Local speedups disappear under composition & Which boundary is hard, which objective is soft, and who wins on conflict? & precedence-aware decision kernel \\
Lifecycle debt accumulates across updates and movement & Where are admit, place, mutate, compact, expose, and evict made explicit? & lifecycle pipeline with bounded actuation interfaces \\
Recovery reuses the same state objects as steady state & Which invariants survive migration, replay, rebuild, or restore? & safety envelope and ownership-transfer rules \\
\bottomrule
\end{tabularx}
\end{table}

\subsection{Cross-Cutting Principles}

Observation before optimization means the relevant state boundary must already be named in telemetry. State as a first-class scheduling object means runtimes should schedule around mutable state rather than passive tasks. Closed-loop rather than open-loop policies means action must feed back into accounting. Stable gains require boundary modeling means every gain must declare what is held fixed, what debt is deferred, and what disturbance invalidates the claim.

\subsection{Design Space Axes}

The architecture axes are state granularity, control-plane coupling, lifecycle management, and service-boundary alignment. Granularity ranges from record to session. Control logic may be embedded, sidecar-based, or shared. Lifecycle management should make admit, place, mutate, compact, and evict explicit. Boundary alignment asks whether local wins survive end-to-end service constraints.

\subsection{Anti-Pattern Repairs}

Telemetry-action mismatch is repaired by naming a legal control action for each signal. Policy layering without precedence is repaired by ordering hard constraints, soft objectives, and tie-breakers. One-shot benchmarking is repaired by adding warmup, disturbance, adaptation, and re-stabilization. Ignoring write amplification is repaired by charging mutation and repair debt to the same service boundary as read gains. Treating ownership as static is repaired by transferable ownership and lease semantics.

\begin{figure*}[t]
\centering
\resizebox{0.96\textwidth}{!}{\begin{tikzpicture}[
  font=\footnotesize,
  ap/.style={
    draw=red!35,
    rounded corners=2pt,
    align=left,
    text width=35mm,
    minimum height=10mm,
    inner sep=3pt,
    fill=red!4
  },
  fix/.style={
    draw=green!35!black,
    rounded corners=2pt,
    align=left,
    text width=37mm,
    minimum height=10mm,
    inner sep=3pt,
    fill=green!5
  },
  mid/.style={
    draw=black!30,
    rounded corners=2pt,
    align=center,
    text=black!62,
    text width=30mm,
    minimum height=11mm,
    inner sep=2.5pt,
    font=\scriptsize,
    fill=black!1
  },
  head/.style={font=\scriptsize\bfseries, align=center, text=black!70},
  bridge/.style={-{Stealth[length=1.7mm]}, line width=0.64pt, draw=black!58, shorten <=2pt, shorten >=2pt},
  spine/.style={line width=0.72pt, draw=black!28}
]

\node[head] at (-1.25,1.98) {Implicit anti-pattern};
\node[head] at (4.25,1.98) {Missing boundary};
\node[head] at (9.75,1.98) {Explicit contract repair};
\draw[black!45, line width=0.6pt] (-3.35,1.80) -- (11.85,1.80);

\node[ap] (a1) at (-1.25,0.95) {\textbf{Telemetry-action mismatch}\\signals do not name legal actions};
\node[ap] (a2) at (-1.25,-0.45) {\textbf{Policy layering}\\objectives lack precedence};
\node[ap] (a3) at (-1.25,-1.85) {\textbf{One-shot benchmark}\\disturbance is omitted};
\node[ap] (a4) at (-1.25,-3.25) {\textbf{Hidden write path}\\read gains defer mutation debt};
\node[ap] (a5) at (-1.25,-4.65) {\textbf{Static ownership}\\handoff legality is implicit};

\node[mid] (m1) at (4.25,0.95) {signal $\rightarrow$ state object};
\node[mid] (m2) at (4.25,-0.45) {ordering and tie-break};
\node[mid] (m3) at (4.25,-1.85) {disturbance phases};
\node[mid] (m4) at (4.25,-3.25) {mutation accounting};
\node[mid] (m5) at (4.25,-4.65) {transfer survival};

\node[fix] (f1) at (9.75,0.95) {\textbf{Typed observation layer}\\metrics feed bounded actions};
\node[fix] (f2) at (9.75,-0.45) {\textbf{Precedence-aware kernel}\\hard rules precede soft goals};
\node[fix] (f3) at (9.75,-1.85) {\textbf{Multi-phase envelope}\\warmup, burst, drift, recovery};
\node[fix] (f4) at (9.75,-3.25) {\textbf{Auditable state catalog}\\writes charged to same boundary};
\node[fix] (f5) at (9.75,-4.65) {\textbf{Transferable state model}\\authority, replay, reclamation};

\foreach \i in {1,2,3,4,5} {
  \draw[bridge] (a\i.east) -- (m\i.west);
  \draw[bridge] (m\i.east) -- (f\i.west);
}
\end{tikzpicture}
}
\caption{Anti-pattern to contract repair map.}
\label{fig:supp-antipattern-map}
\end{figure*}

\subsection{Blueprint Contract Notes}

The blueprint is a conceptual reference model, not a monolithic runtime. Its five layers are typed state catalog, observation and disturbance detection, boundary model and decision kernel, actuation interfaces, and safety and recovery envelope.

\begin{table*}[t]
\small
\caption{Mapping recurring anti-patterns to the blueprint contracts that repair them.}
\label{tab:supp-antipattern-blueprint}
\begin{tabularx}{\textwidth}{>{\raggedright\arraybackslash}p{0.22\textwidth}>{\raggedright\arraybackslash}p{0.26\textwidth}>{\raggedright\arraybackslash}p{0.18\textwidth}>{\raggedright\arraybackslash}Y}
\toprule
Anti-pattern & Missing contract & Primary blueprint layer & Expected design effect \\
\midrule
Telemetry-action mismatch & signals must name a legal actuation target and its delay/confidence semantics & observation and disturbance detection & dashboards become control inputs rather than passive monitoring \\
Policy layering without conflict semantics & hard constraints, soft objectives, and tie-breakers must be ordered explicitly & boundary model and decision kernel & overlapping controllers stop oscillating under burst or skew \\
One-shot benchmarking of non-stationary mechanisms & disturbance phases and re-stabilization windows must be part of the contract & observation plus safety and recovery envelope & reported gains remain meaningful outside stationary replay \\
Ignoring write amplification in read-optimized designs & mutation, compaction, deletion repair, and maintenance debt must be charged to the same boundary as read gains & state catalog plus actuation interfaces & deferred rebuild or compaction cost becomes visible during design review \\
Treating state ownership as static & ownership transfer and no-eviction rules must survive migration and recovery & actuation interfaces plus safety and recovery envelope & scale-out and fault handling share one transferable state model \\
\bottomrule
\end{tabularx}
\end{table*}

\begin{table*}[t]
\small
\caption{Abstract policy skeleton for unified state governance.}
\label{tab:supp-policy-skeleton}
\begin{tabularx}{\textwidth}{>{\raggedright\arraybackslash}p{0.09\textwidth}Y}
\toprule
Step & Policy skeleton \\
\midrule
1 & \texttt{observe\_signals()} and update telemetry confidence for contention, queueing, local freshness, exposure freshness, and retention pressure. \\
2 & \texttt{identify\_state\_objects()} whose boundary is currently active: hot operator shards, KV pages, retrieval indexes, replay buffers, or demotion candidates. \\
3 & \texttt{check\_hard\_constraints()} for SLA, memory limit, recovery envelope, shard-exposure safety, and quality floor before any optimization action. \\
4 & \texttt{rank\_actions()} over admit, place, migrate, refresh, compact, expose, demote, and evict using expected utility and maintenance debt. \\
5 & \texttt{apply\_bounded\_action()} with explicit rollback or retry semantics when telemetry is delayed or conflicting. \\
6 & \texttt{account\_outcome()} by charging the decision against tail latency, quality drift, effective budget loss, and long-horizon reuse value. \\
\bottomrule
\end{tabularx}
\end{table*}

\subsection{Case Study Bridges}

The main paper keeps only bridge-level case text. Full versions should remain here. Multi-tenant LLM serving under burst arrival shows how admission, transfer, reclamation, and restoration share one p99-facing boundary. Stateful stream reconfiguration shows that migration, progress, and replay need one transfer protocol. Retrieval under corpus drift shows that freshness, rebuild debt, and exposure legality must be co-governed.

\subsection{Additional Evidence Anchors for the Hardware-Conscious Cluster}

The hardware-conscious cluster is broader than the main paper can narrate in detail, so the supplement should preserve a few extra anchors that keep the boundary between representation, placement, and movement visible. Approximate execution and compressed-state systems such as Anda, TurboAttention, QServe, ATOM, VAttention, FlashInfer, and SampleAttention make the same point from different angles: representation choice is never free, because it reshapes memory traffic, kernel regularity, and quality debt at the same time~\cite{fang2025anda,kang2025turboattention,lin2025qserve,zhao2024atom,prabhu2025vattention,ye2025flashinfer,zhu2024sampleattention}. Systems such as Chimera, HeterInfer, LIA, HETIS, and Mercury show that heterogeneity-aware placement only pays when the runtime can jointly price compute skew, cache residency, and cross-device movement~\cite{huang2025chimera,chen2025heteroinfer,kim2025lia,mo2025hetis,mercury2025rms}. Others such as WaferLLM, DRAMCache, Tiertune, MoE-lightweight offloading, and PhoenixOS expose how the placement problem shifts once the memory substrate is no longer flat~\cite{he2025waferllm,hong2024dramcache,lee2025tiertune,cao2025moelightning,wei2025phoenixos}.

The same logic extends to storage and I/O-adjacent systems that were intentionally kept out of the main spine. AutoScratch, PFSCK, PolarDB, CXL-based designs, and Toleo all demonstrate that movement, repair, and metadata handling become the dominant cost once the system crosses a hardware boundary that the runtime can no longer ignore~\cite{fu2024autoscratch,domingo2021pfsck,cao2020polardb,ji2025cxl,dong2024toleo}. These are not separate themes from state management; they are the clearest instances of how state representation and transport become a control problem under hardware pressure.

Another useful way to read the remaining hardware-conscious papers is by which part of the boundary they make explicit. Some papers primarily expose representation and quantization debt, including QServe, Atom, QuantLLM, and low-bit cache work~\cite{lin2025qserve,zhao2024atom,xia2024quantllm,dong2024qhitter}. Others expose topology and movement debt, including WaferLLM, JIT-serve, MLP offload, and heterogeneous cache placement lines~\cite{he2025waferllm,zhang2026jitserve,maurya2025mlpoffload,liakopoulos2025maveriq}. A third group exposes fault-tolerance or isolation debt, including Save, ChecknRun, and CXL memory-protection work~\cite{zheng2025save,eisenman2022checknrun,zhao2025equilibria}. Treating those as one cluster preserves the survey's central claim that the runtime must price movement, representation, and resilience together.

The tail of the hardware cluster also includes several papers that are better understood as pressure refinements rather than separate subtopics. FuseMax and related low-overhead quantization systems make the cost of representation smaller, but only by moving complexity into selection and calibration~\cite{nayak2024fusemax,kwon2020fvm,jha2025hycache}. Heterogeneous-storage and memory-tier lines such as Colloid, BeyondHotness, and DRAMCache show that performance depends on where state is allowed to reside and how quickly it can move across tiers~\cite{vuppalapati2024colloid,liu2025beyondhotness,hong2024dramcache}. Specialized accelerator and system-offload papers such as PowerInfer, PIMBA, and QFactory make the same point from the compute side: the runtime only wins when it can retime or relocate state without making control overhead dominate the path~\cite{song2024powerinfer,kim2025pimba,zhang2025qfactory}.

The remaining hardware-concern papers fit the same pattern. MoC and related scheduling lines show that small changes in control policy can shift large amounts of debt between queueing and movement~\cite{cai2025moc}. Fermata-style or optimizer-driven offload systems such as Optimus and Tigon expose where placement and bandwidth interact with the control loop~\cite{feng2025optimus,huang2025tigon}. Finally, system-level memory layout papers such as VM-control, Tabatabai-style memory management, and CXL-aware tiering keep the same point visible: once state moves across a hardware boundary, representation, legality, and cost accounting have to be decided together~\cite{zheng2024vmcu,tabatabai2024fbmm,zhong2024memstrata}.

Finally, several of the remaining boundary papers are best read as support for the same control logic rather than as separate mechanisms. Heterogeneity-aware execution and server placement work such as GMLake, Flex-MoE, and Tigon exposes the interaction between load balance, bandwidth, and local memory pressure~\cite{guo2024gmlake,cao2025moelightning,huang2025tigon}. Approximation and repair papers such as TurboAttention, SampleAttention, and PFSCK show how quality, compute, and metadata repair are coupled once the runtime starts trimming state to fit within the hardware envelope~\cite{kang2025turboattention,zhu2024sampleattention,domingo2021pfsck}. That is the right level of abstraction for the supplement: a boundary map, not an isolated list of systems.

The same repair line also includes concealing and compression-aware execution such as Concealing, which shows that state repair can be hidden inside pipeline slack only when the runtime can separate representation cost from visible service debt~\cite{jin2024concealing}. That reference closes the last gap in the hardware cluster and keeps the supplement honest about the full cost of state movement and repair.

The last few hardware references follow the same pattern and are worth keeping visible in the supplement. Federici-style DP/placement work, Gerogiannis-style decoupled execution, Gu-style centering, and Joo-style coruscant scheduling all show that once the runtime prices movement and placement together, the exact mechanism matters less than the fact that the state boundary is explicit~\cite{federici2025dip,gerogiannis2025deca,gu2025cent,joo2025coruscant}. Kamath-style prefill/decode overlap, Park's Clap and DecDec, and Zico-style low-level control add the same lesson from the scheduling side: overlap only helps if the control overhead stays below the debt it avoids~\cite{kamath2025podattention,park2025clap,park2025decdec,zico21}. Those references are small but useful, because they keep the supplement's hardware map complete without changing the main paper's thesis.

Three more hardware papers round out the same map. Pipeline-style memory optimization systems such as MEPipe, clone-aware movement lines such as Clone, and CXL-aware VM-control work show that the runtime can often trade off recomputation, copying, and placement only because it knows exactly where the state boundary sits~\cite{sun2025mepipe,tian2025clone,zheng2024vmcu}. Memory-heavy systems such as Tigon and MemStrata keep the same lesson visible from the other side of the boundary: when state is split across devices or tiers, legality and placement are still the same control problem~\cite{huang2025tigon,zhong2024memstrata}. Together with the other papers above, they make the hardware cluster complete enough for the supplement without asking the main paper to carry the details.
